\documentclass[11pt,a4paper]{article}
\usepackage[T1]{fontenc}
\usepackage[utf8]{inputenc}
\usepackage{lmodern}
\usepackage[margin=23mm]{geometry}
\usepackage{amsmath,amssymb,booktabs,array,microtype}
\usepackage{xcolor}
\usepackage{hyperref}
\hypersetup{colorlinks=true,urlcolor=blue!55!black,citecolor=blue!55!black,
  linkcolor=blue!55!black,pdftitle={Bounded Lane Detection on Smartphones},
  pdfauthor={YouSpeed project}}
\setlength{\parindent}{0pt}
\setlength{\parskip}{5pt}
\setlength{\emergencystretch}{2em}
\newcommand{\code}[1]{\texttt{\small #1}}
\newcommand{\doc}[1]{\nolinkurl{docs/#1}}
\title{\vspace{-15mm}\textbf{Bounded Lane Detection on Smartphones}\\[3pt]
  \large YouSpeed: methods, engineering contributions, and early evidence}
\author{YouSpeed project\quad\textit{Technical report draft}}
\date{30 September 2026}
\begin{document}
\maketitle
\vspace{-7mm}
\begin{abstract}
\noindent
We describe a classical, on-device lane-boundary pipeline for the YouSpeed
iPhone and Android dashcam previews. Its main contributions are an
output-preserving sparse preparation path, explicitly bounded temporal evidence,
geometry-constrained presentation, and a shared native replay methodology.
A paired Moto benchmark reduced median preparation time from 5.04 to 2.95\,ms;
7,663 recorded frames retained identical tested outputs after optimization.
These are engineering results, not evidence of general lane-detection accuracy.
Lighting experiments and an unsuccessful iPhone field deployment expose remaining
semantic, continuity, and runtime limitations.
\end{abstract}

\section{Objective and scope}
The system estimates visible road-boundary evidence from an existing monocular
camera stream while the phone also supports recording and traffic-sign recognition
(TSR). The optional overlay is a perception prototype, not a lane-departure warning
or vehicle-control system. Missing evidence may produce zero or one displayed
boundary. Confidence is heuristic; neither a plausible pair nor smooth rendering
proves that the boundaries enclose the vehicle's lane.

This report covers the curved, multi-boundary \code{RoadBoundaryDetector} and
\code{RoadPathSession} pipeline, including the September 30 preview candidate.
It supersedes the straight-line description of the initial September prototype.
Results below are attributed to their recorded experiments rather than to one
uniformly qualified release. The reviewed working tree includes uncommitted work
on top of revision \code{a0df46d8c157}; historical measurements were not rerun for
this report.

\section{Relationship to prior work}
Aly's classical method uses inverse perspective mapping, selective oriented
filtering, RANSAC line initialization, and B\'ezier fitting~\cite{aly}.
Our method shares the filter--associate--fit structure, but operates on sparse
image-space rows; its final B\'ezier stage is only a bounded display approximation.
It does not reproduce Aly's detector or inherit its reported performance.

\"Ozgunalp et al. extend symmetrical local thresholding by testing the parallelism
of lane-marking borders and combining this with stereo-derived road segmentation
and bird's-eye geometry~\cite{slt}. Their work motivates our offline border
qualification experiments; our monocular implementation does not provide their
stereo road mask.

Learned alternatives include UFLDv2, which predicts sparse lane coordinates by
hybrid-anchor ordinal classification~\cite{ufld}, and TwinLiteNet+, which jointly
segments lane markings and drivable area~\cite{twin}. These are useful future
comparators, not evaluated baselines in this project. Published benchmark scores
and accelerator timings are not directly comparable with our partial paint labels
or complete phone preparation cost. Classical filtering, patch tracking, and
B\'ezier curves themselves are established techniques; our contribution is their
specific bounded integration and its verification.

\newpage
\section{Method and project-specific contributions}
\subsection{Sparse preparation with equivalent detector inputs}
The camera Y plane is sampled into an upright, aspect-preserving image no larger
than $384\times216$. An immutable sampling plan caches source row and column
offsets for each geometry, including rotation and stride, without retaining the
camera buffer. A five-pixel horizontal morphological opening $I\circ B$ yields
the enhanced image
\begin{equation}
  E=\min\!\left(255,\ I+\left\lfloor\frac{3}{2}
       \bigl(I-(I\circ B)\bigr)\right\rfloor\right).
\end{equation}
Only the detector's support rows are enhanced: at most 24 scan centers, each
reading three neighboring rows, require at most 72 of 216 rows. Raw luma and
enhanced support-row bytes match the full preparation path; unused enhanced rows
are zero. Temporal tracking always reads original luma. This dependency-aware
optimization preserves the detector inputs while avoiding unnecessary work~\cite{prep}.

At each scan row, bilateral bright-stripe contrast identifies paint-like cues;
separate intensity edges remain distinguishable. Local suppression retains at
most 12 candidates per row. Bounded inter-row association builds piecewise
boundaries rather than imposing a globally straight lane. Only suitable paint
evidence forms corridor hypotheses; an edge alone cannot establish one.
Visual calibration chooses the scan horizon and gives nearby observed candidates
a bounded ranking preference, without creating observations.

\subsection{Evidence-aware temporal continuity and preview selection}
The tracker searches original-luma patches with forward/backward agreement and
coherent residual checks. Boundaries record fresh, tracked, or fused provenance
and evidence age. Eligible calibration, causal speed, and GNSS heading rate can
predict approximate road-plane motion to center the bounded patch search.
Image correspondence must still support the prediction. This is not a complete
world-coordinate estimator or an exposure-synchronized IMU solution~\cite{stability}.

The preview runs independently of TSR inference, targeting one admission every
100\,ms with one active and one latest pending image. Preparation has a 50\,ms
budget and operation caps; scope changes and thermal pauses invalidate work.
The target rate is not a sustained-throughput claim. After maturity filtering,
the selector publishes at most one boundary on each side, checks pair ordering
and width, and requires a sustained 0.35\,s advantage before replacing a visible
incumbent. Identity can survive a short absence up to 0.75\,s, but absent
coordinates are not republished. Calibration guides remain visually distinct.
The preview and TSR sidecar currently prepare geometry separately, adding cost.

\subsection{Bounded rendering and native reproducibility}
Display polylines are approximated by cubic B\'ezier segments satisfying
\begin{equation}
 \max_{y\ \text{in each supported interval}}
   |x_{\mathrm{Bezier}}(y)-x_{\mathrm{polyline}}(y)|\leq 2/383.
\end{equation}
This is a horizontal error bound of two pixels at the maximum analysis width.
Endpoints are retained; invalid points, non-increasing $y$, and large gaps split
curves. Rendering does not extrapolate beyond support or change detection,
confidence, or TSR evidence~\cite{bezier}.

Swift and Kotlin implement matching arithmetic and contracts. Shared fixtures,
native video replay, input/source hashes, and separate raw-versus-published
geometry make algorithm changes auditable. This combination of bounded evidence,
output-preserving optimization, and parity checks constitutes our engineering
innovation; priority over all prior systems is not claimed.

\newpage
\section{Evaluation and limitations}
\textbf{Preparation cost and equivalence.}
The paired Moto g86 5G experiment alternated baseline and optimized execution
within one standalone ART process for 60 seconds, with 7,149 measured samples per
variant after warm-up. It repeated three fixtures, including live $288\times216$
analysis geometry. Preparation included sampling, filtering, detection, tracking,
and presentation preparation; it excluded camera delivery, TSR, recording, and
UI~\cite{prep}.

\begin{center}
\small
\begin{tabular}{@{}lrr@{}}
\toprule
Measured warm preparation & Baseline & Optimized \\
\midrule
Median wall time & 5.04\,ms & 2.95\,ms \\
95th-percentile wall time & 6.14\,ms & 3.88\,ms \\
Median thread CPU time & 4.99\,ms & 2.92\,ms \\
50\,ms budget failures & 0/7,149 & 0/7,149 \\
\bottomrule
\end{tabular}
\end{center}
Median and p95 wall time fell by 41.5\% and 36.8\%, respectively. Separately,
all 7,663 replayed frames matched the pre-change tested outputs exactly;
26 boundary and 21 path fixtures retained Swift/Kotlin parity within $10^{-9}$.
These support local computation savings and regression equivalence, not improved
accuracy or sustained full-app thermal performance. B\'ezier parity covered 2,637
boundaries, reducing 28,090 line segments to 6,293 cubic segments~\cite{bezier}.

\textbf{Continuity candidate.}
Three development clips supplied 600 exposures over 59.73\,s at approximately
10\,Hz. Compared with the build 10022 core, the final preview candidate reduced
visible identities from 97 to 49 and blank frames from 172 to 148; frames with
more than two displayed lines fell from 42 to zero. Identity-set transitions
remained 216. Raw detector/tracker output was unchanged because replay supplied
no motion calibration. Thus these results measure selection and maturity effects,
not improved accuracy or validated motion compensation~\cite{stability}.

\textbf{Lighting experiments.}
Eight frozen variants processed 7,663 frames from two recordings, sampled at
2\,Hz. Evaluation used only 58 partially annotated frames, including six short
new sequences. At three-pixel tolerance, local normalization increased newly
labelled visible-paint coverage from 733/2,130 (34.4\%) to 817/2,130 (38.4\%),
but lost correct paint and produced false roadside geometry in visual review.
Border qualification showed no primary visible gain. No variant was promoted.
These reused drives are not an untouched holdout; coverage is neither precision
nor full-lane accuracy, and unlabelled pixels cannot be treated as negatives~\cite{lighting}.

\textbf{Field failure and next validation.}
The iPhone 14 Pro build 10023 recorded 744/744 preview preparation-budget failures
and no visible boundaries, followed by thermal pause. The deployed Swift Debug
build used \code{-Onone}; a matched 150-frame Mac replay failed on all frames
without optimization and on none with \code{-O}. This isolates a plausible runtime
mechanism, but does not establish optimized iPhone performance. A separate probe
confirmed that tiny intrinsic-value changes reset tracking identity, independently
of preparation cost~\cite{failure}.

The next acceptance gate is optimized full-app testing on both phones, stable
calibration identity, and an untouched drive with consecutive labels. It should
measure false-visible duration, correct-side availability, reacquisition,
capture-to-overlay age, TSR contention, recording continuity, and thermal behavior.
Encoded-video crops and live camera geometry must be checked separately.
The demonstrated advance is a cheaper, reproducible mobile pipeline; robust
on-road lane interpretation remains an open validation task.

\newpage
\begin{thebibliography}{10}
\small
\setlength{\itemsep}{7pt}
\bibitem{aly}
M. Aly, ``Real time detection of lane markers in urban streets,''
\textit{IEEE Intelligent Vehicles Symposium}, 2008.
\href{https://doi.org/10.1109/IVS.2008.4621152}{doi:10.1109/IVS.2008.4621152}.
\href{https://arxiv.org/abs/1411.7113}{Open manuscript: arXiv:1411.7113}.

\bibitem{slt}
U. \"Ozgunalp et al., ``Robust lane marking detection algorithm using drivable
area segmentation and extended SLT,'' \textit{IEEE ROBIO}, 2019.
\href{https://arxiv.org/abs/1911.09054}{arXiv:1911.09054}.

\bibitem{ufld}
Z. Qin, P. Zhang, and X. Li, ``Ultra fast deep lane detection with hybrid anchor
driven ordinal classification,'' \textit{IEEE TPAMI}, 2022, early publication.
\href{https://arxiv.org/abs/2206.07389}{arXiv:2206.07389}.

\bibitem{twin}
Q.-H. Che, D.-T. Le, M.-Q. Pham, V.-T. Nguyen, and D.-K. Lam,
``TwinLiteNet+: An enhanced multi-task segmentation model for autonomous driving,''
\textit{Computers and Electrical Engineering}, vol. 128, article 110694, 2025.
\href{https://arxiv.org/abs/2403.16958v6}{arXiv:2403.16958v6, revised 2026}.

\bibitem{prep}
YouSpeed project, ``Lane preparation optimization,'' internal experiment report,
30 September 2026. \doc{LANE_PREPARATION_OPTIMIZATION_2026-09-30.md}.
Includes exact replay comparisons, sampler/filter tests, the paired Moto
benchmark, and a rejected whole-track verification experiment.

\bibitem{stability}
YouSpeed project, ``Lane stability implementation,'' development-candidate report,
30 September 2026. \doc{LANE_STABILITY_IMPLEMENTATION_2026-09-30.md}.
Includes the 600-frame paired preview comparison and rejected variants.

\bibitem{bezier}
YouSpeed project, ``Lane presentation and calibration review,''
30 September 2026. \doc{LANE_BEZIER_AND_CALIBRATION_REVIEW_2026-09-30.md}.
Includes the display error bound and recorded-boundary native parity check.

\bibitem{lighting}
YouSpeed project, ``Lane lighting evaluation,'' offline experimental report,
30 September 2026. \doc{LANE_LIGHTING_EVALUATION_2026-09-30.md}.
Includes frozen variants, partial-label methodology, tolerance sensitivity,
and qualitative false-boundary review.

\bibitem{failure}
YouSpeed project, ``iPhone build 10023: lane failure diagnosis,''
30 September 2026. \doc{LANE_IPHONE_BUILD10023_FAILURE_2026-09-30.md}.
Includes device diagnostics, matched optimization replay, and calibration-churn probe.
\end{thebibliography}

\vspace{5pt}
{\small\textbf{Reproducibility note.}
Internal paths are relative to the YouSpeed repository. Source entry points are
\nolinkurl{iphone/SpeedConsumerApp/} and
\nolinkurl{android/app/src/main/java/de/youspeed/android/alpha/}, principally
\code{RoadBoundaryDetector}, \code{RoadBoundaryTemporalTracker},
\code{RoadBoundaryPresentationGate}, \code{RoadPathSession}, and
\code{LaneBoundaryBezier}. Replay tools are under \nolinkurl{scripts/lanes/}.
The cited reports identify frozen sources and hashed outputs. Driving recordings
and detailed evidence are retained locally in ignored directories and are not
distributed with this standalone document.}
\end{document}
